\section{Identifiability and the singleton device}
\label{sec:identifiability}

The second result the siblings share is structural identifiability:
when do the coarsened reports determine $\theta$, and what restores
identifiability when they do not. We state the rank condition once and
the singleton-restoration result once, then tabulate the per-domain
singleton devices as instances of one mechanism.

\subsection{The augmented-candidate-set rank condition}

Let $\mathcal{S}$ be the collection of candidate sets that occur with
positive probability under the coarsening, and let $C \in
\{0,1\}^{|\mathcal{S}| \times m}$ be the candidate-set matrix whose
rows are the indicator vectors of the sets in $\mathcal{S}$ over an
$m$-element latent support. Write $\tilde C$ for $C$ augmented with the
all-ones row, the row contributed by the fully uninformative report.

\begin{theorem}[Identifiability under C1--C2--C3]
\label{thm:general-rank}
Suppose each latent component's contribution to the report law is
individually identifiable, and C1--C2--C3 (\cref{cond:c1,cond:c2,cond:c3})
hold. Then:
\begin{enumerate}[(a)]
\item \textbf{Necessity.} If two latent components are
  \emph{coarsening-confounded}, that is, every candidate set in
  $\mathcal{S}$ contains both or neither, their parameters are not
  separately identifiable from the coarsened-data likelihood: any
  reparametrization preserving their summed contribution leaves the
  likelihood unchanged.
\item \textbf{Sufficiency.} If $\tilde C$ has full column rank $m$, then
  $\theta$ is identifiable.
\end{enumerate}
\end{theorem}

\begin{remark}[Instantiating the rank condition]
\label{rem:rank-instantiation}
The sufficiency hypothesis is the column-rank condition on $\tilde C$
alone; each domain then supplies whatever support condition makes its
own coarsening map produce a full-rank incidence. In the time-to-event
instance this is the requirement that the mechanism assign positive
probability to both exact and censored reports, so that the all-ones
augmenting row is realized; the continuous-report instances replace it
with the analogous support condition on the operator
(\cref{thm:general-rank}, proof sketch). These are instantiation
details of one rank condition, not separate hypotheses of the general
theorem.
\end{remark}

\begin{proof}[Sketch]
This is Theorem on identifiability under C1--C2--C3 of
\citet{towell2026masked}, restated for the abstract process. The
necessity direction is the observation that under C1--C2--C3 the
likelihood depends on a confounded pair only through the sum of their
contributions, so the sum is identifiable but the split is not. The
sufficiency direction differences the report law across candidate sets
to obtain $\tilde C\, g(\cdot) = 0$ for the vector $g$ of pairwise
contribution differences, whence full column rank of $\tilde C$ forces
$g = 0$ and individual identifiability closes the argument. The
continuous-report instances replace the finite incidence matrix
$\tilde C$ with its operator analogue: in spatial transcriptomics it
is the spot-by-cell-type composition matrix $P$, identifiable exactly
when $P$ has full column rank
\citep{towell2026spatialcoarsening}; in differential privacy with a
linear query $q(D) = AD$ it is the design matrix $A$, identifiable
exactly when $A$ has full column rank \citep{towell2026dpcoarsening};
in weak supervision it is the agreement matrix whose off-diagonal is
rank one under conditional independence, with a rank deficit $r$
induced by labeling-function dependence
\citep{towell2026weaksupcoarsening}; in multiple instance learning it
is the bag-by-instance-type composition matrix $M$, identifiable exactly
when $M$ has full column rank \citep{towell2026milcoarsening}. Each is
the same rank condition
on the augmented coarsening operator. The autopsy-model ancestry of
the finite case is \citet{meilijson1981}; the competing-risks
nonidentifiability without independence is \citet{tsiatis1975}.
\end{proof}

\subsection{Singletons restore identifiability}

When the rank condition fails, every domain reaches for the same
remedy: a report that pins the latent value to a point.

\begin{proposition}[Singleton restoration]
\label{prop:singleton}
A report $r$ with $|c(r)| = 1$, a \emph{singleton} candidate set,
contributes the latent value exactly. Augmenting the candidate-set
collection $\mathcal{S}$ with singletons for components on which
$\tilde C$ is rank-deficient raises the column rank of $\tilde C$; a
collection of singletons that covers a basis of the deficient subspace
restores full column rank and hence identifiability under
\cref{thm:general-rank}. A singleton satisfies C2 vacuously, so it adds
information without adding a modeling burden.
\end{proposition}

\begin{proof}[Sketch]
A singleton $\{j\}$ adds the standard basis row $e_j^\top$ to $C$.
Adding rows can only raise rank, and a set of singletons spanning the
left-null directions of $\tilde C$ removes the rank deficit, at which
point \cref{thm:general-rank}(b) applies. The C2-vacuousness of
singletons is \cref{cond:c2} at $|c(r)| = 1$. The quantitative version,
how many singletons are needed, is domain-specific: in weak
supervision it is set by the rank deficit $r$ and the labeling-function
accuracy margin \citep{towell2026weaksupcoarsening}; in the others a
basis-covering set suffices.
\end{proof}

The content of \cref{prop:singleton} is that the recovery devices the
six fields invented independently are one device. \Cref{tab:singletons}
makes the identification explicit.

\begin{table}[tb]
\centering
\small
\begin{tabular}{@{}p{0.20\linewidth} p{0.34\linewidth} p{0.38\linewidth}@{}}
\toprule
\textbf{Domain} & \textbf{Singleton device} & \textbf{What it pins} \\
\midrule
Reliability \citep{towell2026masked} &
  A diagnostic that resolves the failed component &
  The exact cause $K_i$ \\
scRNA-seq \citep{towell2026scrnacoarsening} &
  ERCC spike-in transcript (known input) &
  The dropout function, hence the per-gene mean \\
Spatial \citep{towell2026spatialcoarsening} &
  Single-cell-resolution probe / spot &
  The cell-type-specific expression at that location \\
Differential privacy \citep{towell2026dpcoarsening} &
  A non-private (noise-free) release &
  The kernel scale, hence the functional $\theta$ \\
Weak supervision \citep{towell2026weaksupcoarsening} &
  A gold-labeled example &
  The true label $Y_i$ \\
Phenotyping \citep{towell2026phenotypecoarsening} &
  A chart-reviewed (adjudicated) patient &
  The true disease state, hence sens/spec \\
Multiple instance learning \citep{towell2026milcoarsening} &
  A singleton bag (an instance-level label) &
  The instance's true positivity, hence instance-level scores \\
\bottomrule
\end{tabular}
\caption{The singleton candidate set in each domain. These are not
analogies: each is a report $r$ with $|c(r)| = 1$ in the sense of
\eqref{eq:fw-candidate}, and each restores identifiability by the same
rank argument (\cref{prop:singleton}).}
\label{tab:singletons}
\end{table}

\subsection{When C2 fails}

The conditions are not always met, and the framework says what breaks.
When C2 fails, the coarsening weight no longer factors out, the score
acquires an unmodeled $\theta$-dependent term, and the affected
functionals lose identifiability from coarsened reports alone. This is
the same failure in every domain: informative masking in reliability
\citep{towell2026mdrelax}, data-dependent mechanisms (sparse vector,
propose-test-release) in differential privacy
\citep{towell2026dpcoarsening}, labeling-function dependence in weak
supervision \citep{towell2026weaksupcoarsening}, and verification bias
on a non-random chart-reviewed subsample in phenotyping
\citep{towell2026phenotypecoarsening}. The remedy is again the
singleton of \cref{prop:singleton}, plus, where available, the
sensitivity bands of \citet{towell2026mdrelax} that quantify how far
estimates can move under a bounded C2 violation.
